\documentclass[11pt,a4paper]{article}
\usepackage[margin=1in]{geometry}
\usepackage{amsmath,amssymb,amsthm}
\usepackage{algorithm}
\usepackage{algpseudocode}
\usepackage{booktabs}
\usepackage{hyperref}

\theoremstyle{definition}
\newtheorem{definition}{Definition}[section]
\newtheorem{theorem}{Theorem}[section]
\newtheorem{lemma}[theorem]{Lemma}
\newtheorem{remark}{Remark}[section]

\title{\textbf{Rigorous Mathematical Specification and Multi-Horizon Forecast of Multi-Token Prediction}}
\author{Team Scaibu \\ \small Email: \texttt{jaydeep.9664920749@gmail.com} \\ \small Architecture and Core Engine Team}
\date{October 2026}

\begin{document}
\maketitle

\begin{abstract}
Standard language models optimize next-token prediction $\log P(x_{t+1} \mid x_{\le t})$. Multi-Token Prediction trains $K$ independent output heads branching from a shared transformer trunk to simultaneously forecast $K$ future tokens $(x_{t+1}, \dots, x_{t+K})$. This specification formalizes the multi-head loss objective, proving joint likelihood optimization.
\end{abstract}

\begin{equation}
\mathcal{L}_{\text{MTP}} = \frac{1}{K} \sum_{k=1}^K -\log P_k(x_{t+k} \mid \mathbf{h}_t)
\end{equation}

\begin{thebibliography}{9}
\bibitem{gloeckle2024} F.~Gloeckle et al., ``Better & faster large language models via multi-token prediction,'' \emph{arXiv:2404.19737}, 2024.
\end{thebibliography}

\end{document}
